\documentclass[10pt,a4paper]{amsart}
\usepackage[margin=19mm]{geometry}
\usepackage{amsmath,amssymb,mathtools}
\usepackage[scr=rsfs,cal=euler]{mathalfa}
\usepackage{xfrac}
\usepackage[unicode,hidelinks,bookmarks=false]{hyperref}
\newcommand{\ie}{i.e.\ }
\newcommand{\cf}{cf.\ }
\newcommand{\resp}{respectively\ }
\newcommand{\Subsection}{\subsection}
\providecommand{\nobreakdash}{\nobreak\hbox{-}}
\setlength{\emergencystretch}{2em}
\multlinegap0pt
\usepackage{bm}
\usepackage{bbm}
\usepackage{dsfont}
\usepackage{xspace}
\usepackage{cancel}
\usepackage{mathtools}
\usepackage{amsthm}
\makeatletter
\@ifundefined{lemma}{\newtheorem{lemma}{Lemma}}{}
\makeatother
\title[Order automorphisms of effect algebras]{Order automorphisms of effect algebras}
\author{Peter Semrl}
\date{Source: arXiv:2602.20888v1 (CC BY 4.0)}
\begin{document}
\maketitle

\noindent\emph{Source and licence.} Order automorphisms of effect algebras. Version arXiv:2602.20888v1 (CC BY 4.0), distributed under the Creative Commons Attribution 4.0 International licence, as recorded in the arXiv licence metadata for this exact version and shown on the abstract page https://arxiv.org/abs/2602.20888v1. Full rights notice: licenses/arxiv-2602.20888-CC-BY-4.0.txt.\par\smallskip

\noindent\textit{Editorial scope.} Counted unit: the Lemma whose source label is \texttt{max} in the article (source0.tex lines 388--398, source label \texttt{max}), delivered together with the article's own complete original proof of it (source0.tex lines 400--453, one \texttt{proof} environment of 54 lines). No mathematics is written, repaired or re-derived: statement and proof are the article's own text line for line. Required context retained uncounted: none. Every definition, display and label of those lines is retained unchanged. Omitted from this package and not redistributed: 1--387, 399--399, 454--1217. Nothing inside a retained excerpt is omitted or altered: every retained body is delivered exactly as the article's lines are written, so no omission receipt of any kind accompanies this package. Citations: the retained excerpts carry no citation command at all, so no bibliography entry of the article is transcribed and no reference is redistributed. Reference adaptation: none was needed and none was made; every reference inside the delivered unit resolves inside it, so no unresolved reference or citation is produced. Printed slips kept literal: no printed word, symbol, hypothesis or number of the counted statement or of its proof was altered. Layout and class: the article's own class is \texttt{article} with packages including amsthm, amssymb, amsfonts, color, amsmath; the wrapper is the lane's pinned \texttt{amsart} editorial wrapper at 10pt on A4, which restates the theorem-like environments the retained text needs and adds the article's own macro definitions that the retained text needs (none), with extra wrapper packages (bm, bbm, dsfont, xspace, cancel, mathtools). The delivered numbering is the unit's own. Assets: no retained span contains a figure environment, a graphics-inclusion command, a PDF-page inclusion, a table or a TikZ picture; no image file is copied into this package, the package declares no asset dependency and no image file is redistributed with it. Licence: version arXiv:2602.20888v1 of this article is distributed under the Creative Commons Attribution 4.0 International licence, as recorded in the arXiv licence metadata for this exact version and shown on the abstract page https://arxiv.org/abs/2602.20888v1. A full rights notice is delivered at licenses/arxiv-2602.20888-CC-BY-4.0.txt. Characters: every retained excerpt is pure ASCII, so the delivered engine, XeLaTeX with its default Latin Modern text font, reports no missing character.
\begin{lemma}\label{max}
Let $A,B \in [0,I] \subset S_2$. The following are equivalent.
\begin{itemize}
\item $\mathcal{D}_A = \mathcal{D}_B$.
\item $A= B$ or $A,B \in \mathcal{R} \cup \{ 0 \}$ or $A = JBJ$,
where
$$
J = \left[ \begin{matrix} 1& 0 \cr 0 & -1 \cr \end{matrix} \right] .
$$
\end{itemize}
\end{lemma}

\begin{proof}
We first note that a $2 \times 2$ symmetric matrix is positive if and only if both diagonal entries are nonengative and its determinant is nonnegative. Indeed, assume first $A \in S_2$ is positive. Let $\{ e_1 , e_2 \}$ be the standard basis of $\mathbb{R}^2$. Then $\langle Ae_1 , e_1 \rangle \ge 0$ and  $\langle Ae_2 , e_2 \rangle \ge 0$ and this tells us that the diagonal entries of $A$ are nonnegative. Since both eigenvalues of $A$ are nonnegative we conclude that $\det A \ge 0$. Assume next that  both diagonal entries of $A$ are nonnegative and $\det A \ge 0$. From the nonnegativity of the determinant we infer that both eigenvalues of $A$ are nonnegative or both are nonpositive and at least one is negative. In the second case the trace of $A$ would be negative contradicting our assumption. Thus, both eigenvalues are nonegative and therefore $A \ge 0$, as desired.

In the next step we will describe the set of maximal elements of $\mathcal{D}_A$ for any $A\in [0,I]$. Let us start with the simplest case that $A \in \mathcal{D}$. Then obviously, the set of maximal elements of $\mathcal{D}_A$ is the singleton $\{ A \}$. Thus we assume from now on that the off-diagonal entries of $A$ are nonzero,
$$
A = \left[ \begin{matrix} t& u \cr u & s \cr \end{matrix} \right] 
$$
with $u \not= 0$. Since $A \ge 0$ we have $ts - u^2 \ge 0$. By the first paragraph of the proof we know that for a pair of real numbers $p,q \in [0,1]$ we have
$$
\left[ \begin{matrix} p& 0 \cr 0 & q \cr \end{matrix} \right] \le \left[ \begin{matrix} t& u \cr u & s \cr \end{matrix} \right]
$$
if and only if $0 \le p \le t$ and $0 \le q \le s$ and 
\begin{equation}\label{detpos}
(t-p)(s-q) - u^2 \ge 0.
\end{equation}
The first case that we will treat is that $A$ is of rank one and has nonzero off-diagonal entries. Then $ts = u^2 \not= 0$. If at least one of the nonnegative numbers $p,q$, $p \le t$, $q \le s$, is positive then
$$
(t-p)(s-q) - u^2 < ts - u^2 = 0 .
$$
It follows that  $\mathcal{D}_A = \{ 0 \}$. 

It remains to consider the case that 
$$A =  \left[ \begin{matrix} t& u \cr u & s \cr \end{matrix} \right]  \in [0,I]
$$ 
is not diagonal and is of rank two. We claim that the set $\mathcal{S}$ of all maximal elements of
$\mathcal{D}_A$ is equal to
$$
\mathcal{S} = \left\{    \left[ \begin{matrix} p& 0 \cr 0 & q \cr \end{matrix} \right] \, : \, 0\le p \le t \ \, {\rm and} \ \, 0 \le q \le s \ \, {\rm and} \ \, (t-p)(s-q) = u^2 
\right\}. 
$$
Let ${\rm diag}\, (p,q) \in \mathcal{S}$. Then clearly ${\rm diag}\, (p,q) \le A$. Take any ${\rm diag}\, (p',q') \in \mathcal{D}$ such that ${\rm diag}\, (p',q') \ge {\rm diag}\, (p,q)$ and ${\rm diag}\, (p',q') \not= {\rm diag}\, (p,q)$. In order to prove that ${\rm diag}\, (p,q)$ is a maximal element of $\mathcal{D}_A$ we need to verify that \begin{equation}\label{mmm}{\rm diag}\, (p',q') \not\le A.\end{equation} If $p' > t$ or $q' > s$ then the verification of (\ref{mmm}) is trivial. So, we may assume that $p' \le t$ and $q' \le s$. We have $p' >p$ or $q' > q$ and therefore $(t-p')(s-q') < (t-p)(s-q) = u^2$ yielding that the determinant of $A - {\rm diag}\, (p',q')$ is negative, and consequently, (\ref{mmm}) holds in this case, as well. Hence, every element of $\mathcal{S}$ is a maximal element of $\mathcal{D}_A$.

Consider now ${\rm diag}\, (p,q) \in \mathcal{D}_A$ such that  ${\rm diag}\, (p,q) \not\in \mathcal{S}$. Then $(t-p)(s-q) > u^2$ yielding that $p <t$ and $q<s$. Then we can find $p'$, $p < p' < t$ such that 
$(t-p')(s-q) > u^2$ which implies that ${\rm diag}\, (p',q) \in \mathcal{D}_A$. Thus, ${\rm diag}\, (p,q)$ is not a maximal element of  $\mathcal{D}_A$. The claim that the set of maximal elements of $\mathcal{D}_A$  coincides with $\mathcal{S}$ has been proved.

We are now ready to prove our equivalence. Assume first that $A$ and $B$ satisfy the second condition. It is clear that if $A=B$ then 
$\mathcal{D}_A = \mathcal{D}_B$. If $A,B \in \mathcal{R} \cup \{ 0 \}$ then 
$\mathcal{D}_A = \mathcal{D}_B = \{ 0 \}$. So, assume that $A= JBJ$. Then obviously $B = JAJ$. If $X \in \mathcal{D}_A$ then $X \le A$ and therefore $X = JXJ  \le JAJ = B$. The proof in one direction is completed.

Hence, assume that $\mathcal{D}_A = \mathcal{D}_B$. We start with the possibility that  $\mathcal{D}_A = \mathcal{D}_B = \{ 0 \}$. Then we know that 
$A,B \in \mathcal{R} \cup \{ 0 \}$. Next, if the set of all maximal elements of $\mathcal{D}_A$ coincides with the set of all maximal elements of $\mathcal{D}_B$ and this set is a singleton $\{ C \}$, $C\not= 0$, then we have already shown that $A=B=C$ is a diagonal matrix. It remains to consider the case when the set of all maximal elements of
$\mathcal{D}_A$, which is equal to  the set of all maximal elements of
$\mathcal{D}_B$, is an infinite set. Then both $A$ and $B$ are invertible matrices with nonzero off-diagonal entries,
$$
A =  \left[ \begin{matrix} t& u \cr u & s \cr \end{matrix} \right] \ \ \ {\rm and} \ \ \ B =  \left[ \begin{matrix} t'& u' \cr u' & s' \cr \end{matrix} \right].
$$
The set of maximal elements of $\mathcal{D}_A$ is the same as the set 
of maximal elements of $\mathcal{D}_B$, that is, the two sets below
$$
\{ (p,q)\, : \, 0\le p \le t \ \, {\rm and} \ \, 0 \le q \le s \ \, {\rm and} \ \, (t-p)(s-q) = u^2  \}$$ and $$
\{ (p,q)\, : \, 0\le p \le t' \ \, {\rm and} \ \, 0 \le q \le s' \ \, {\rm and} \ \, (t'-p)(s'-q) = u'^2  \} 
$$
are the same. From here one can easily conclude that $t=t'$, $s=s'$ and $u = \pm u'$. It follows that $A=B$ or $A = JBJ$.
\end{proof}

\end{document}
